% !TEX program = lualatex
% =====================================================================
%  練習会 最小手順シート（A4 表裏 1 枚・当日持ち歩き用）
%  ・詳しい版：practice.tex（練習会ミニマニュアル・11 ページ）
%  ・本編　　：main.tex（計算センター運営テキストブック）
% =====================================================================
\documentclass[a4paper,9pt]{ltjsarticle}

\usepackage{geometry}
\geometry{top=12mm,bottom=12mm,left=12mm,right=12mm,headsep=3mm,footskip=8mm}

\usepackage{array,longtable,booktabs,colortbl,xcolor,needspace,enumitem,fancyhdr}
\usepackage[most]{tcolorbox}
\usepackage[hidelinks]{hyperref}

\definecolor{cSI}{HTML}{1F5FA8}
\definecolor{cEMIT}{HTML}{9C5A00}
\definecolor{cSIAC}{HTML}{0E7A5F}
\definecolor{j2mbg}{HTML}{EDF0FB}
\definecolor{j2mfr}{HTML}{3A49A8}
\definecolor{notebg}{HTML}{F4F4F0}
\definecolor{notefr}{HTML}{888880}
\definecolor{failbg}{HTML}{FCEEEA}
\definecolor{failfr}{HTML}{B5502E}
\definecolor{band}{gray}{0.90}
\definecolor{rulec}{gray}{0.35}
\definecolor{boxln}{HTML}{444444}

\pagestyle{fancy}
\fancyhf{}
\fancyhead[L]{\small 練習会 最小手順シート}
\fancyhead[R]{\small 計算センター}
\fancyfoot[C]{\small\thepage}
\setlength{\headheight}{14pt}
\renewcommand{\headrulewidth}{0.4pt}

\setlength{\parindent}{0pt}
\setlength{\parskip}{3pt}
\setlength{\fboxsep}{2pt}

\DeclareRobustCommand{\tg}[2]{{\colorbox{#1}{\color{white}\sffamily\bfseries\scriptsize\strut #2}}}
\DeclareRobustCommand{\tSI}{\tg{cSI}{SI}}
\DeclareRobustCommand{\tEMIT}{\tg{cEMIT}{EMIT}}
\DeclareRobustCommand{\tSIAC}{\tg{cSIAC}{SIAC}}

\newcommand{\hd}[1]{\textbf{\sffamily #1}}
\newcolumntype{P}[1]{>{\raggedright\arraybackslash}p{#1}}

\newcommand{\Head}[1]{%
  \Needspace*{4\baselineskip}%
  \par\vspace{5pt}%
  \noindent\colorbox{boxln}{\color{white}\sffamily\bfseries\strut\ #1\ }\par
  \vspace{1pt}\noindent\textcolor{rulec}{\rule{\textwidth}{0.7pt}}\par\vspace{2pt}%
}

\newtcolorbox{Note}[1][]{enhanced,breakable,colback=notebg,colframe=notefr,
  boxrule=0.4pt,left=5pt,right=5pt,top=7pt,bottom=3pt,arc=2pt,
  title={\sffamily\bfseries\footnotesize #1},
  fonttitle=\normalfont,coltitle=white,attach boxed title to top left={xshift=6pt,yshift=-2pt},
  boxed title style={colback=notefr,arc=1pt,boxrule=0pt}}
\newtcolorbox{J2M}[1][JOY2Mulka を使う場合]{enhanced,breakable,
  colback=j2mbg,colframe=j2mfr,
  boxrule=0.5pt,left=5pt,right=5pt,top=7pt,bottom=3pt,arc=2pt,
  title={\sffamily\bfseries\footnotesize #1},
  fonttitle=\normalfont,coltitle=white,attach boxed title to top left={xshift=6pt,yshift=-2pt},
  boxed title style={colback=j2mfr,arc=1pt,boxrule=0pt}}
\newtcolorbox{Fail}[1][つまずきやすいところ]{enhanced,breakable,
  colback=failbg,colframe=failfr,
  boxrule=0.5pt,left=5pt,right=5pt,top=7pt,bottom=3pt,arc=2pt,
  title={\sffamily\bfseries\footnotesize #1},
  fonttitle=\normalfont,coltitle=white,attach boxed title to top left={xshift=6pt,yshift=-2pt},
  boxed title style={colback=failfr,arc=1pt,boxrule=0pt}}

\begin{document}
\gtfamily\sffamily

\begin{center}
{\LARGE\bfseries 練習会 最小手順シート}\quad
{\small フリースタート・PC 1 台・当日申込中心}\\[1mm]
\textcolor{rulec}{\rule{\textwidth}{0.8pt}}
\end{center}

\textbf{このシートは A4 表裏 1 枚です。当日はこれだけ持てば回ります。}
詳しい説明は『練習会の計セン ミニマニュアル』（\texttt{practice.pdf}），
大会運営は本編『計算センター運営テキストブック』（\texttt{main.pdf}）。

\Head{計センがやること（3 つだけ）}

\begin{enumerate}[leftmargin=*,itemsep=0pt]
\item 走る人と\textbf{カード番号}をひもづける
\item コースの\textbf{正解の並び}を Mulka2 に入れておく
\item 帰ってきた\textbf{カードを読む}
\end{enumerate}

記録の計算は Mulka2 がやります。トラブルの大半は，この 3 つのどれかがズレているだけです。

\Head{持っていくもの}

\begin{center}
\small
\begin{tabular}{@{}P{26mm}P{150mm}@{}}
\toprule
\hd{機材} & PC 1 台／読み取り機 1 台\textbf{＋予備 1 台}／USB ケーブル\textbf{＋予備}／
  コントロール用の機材／クリア・チェック・スタート・フィニッシュ \\
\hd{紙} & \textbf{カード番号と氏名の対応表}（カード番号順）／\textbf{救護用リスト}（全員・1 枚）／
  全ポ図・コース図／白紙と筆記具 \\
\hd{電源} & PC の電源とモバイルバッテリー。屋外なら発電機かポータブル電源 \\
\bottomrule
\end{tabular}
\end{center}

\textbf{プリンタは要りません。}結果はあとからクラウド／ラップセンターで公開します。
紙は前日までに刷っておきます。

\Head{前日までにやること}

\begin{enumerate}[leftmargin=*,itemsep=1pt]
\item \textbf{イベントを作る}：Mulka2 のイベントマネージャで新規作成。大会名と日付を入れる。
\item \textbf{コースデータを入れる}：OCAD の \texttt{〜CoursesV8.txt} をドロップ枠に入れる。
      無ければ表計算ソフトで \texttt{Course.csv} を作る
      （列：コース，距離，登距離，コントロール数，S，1，2，…，F）。
\item \textbf{クラスデータを入れる}：\texttt{Class.csv}（クラス，コース，入賞人数，競技時間）。
      \textbf{競技時間は入れる}。空欄だと何時間かかっても完走扱いになる。
\item \textbf{参加者を入れる}：\texttt{Startlist.csv} をドロップ枠へ。
      スタート時刻の列は\textbf{空でよい}（フリースタート）。
\item \textbf{紙を 2 枚刷る}：カード対応表（カード番号順）と救護用リスト（全員）。
\item \textbf{時計を合わせる}：PC の時計。「設定」→【今すぐ同期】が失敗したら，
      管理者権限のコマンドプロンプトで\ \texttt{w32tm /resync /force}。
      確認は\ \url{https://www.nict.go.jp/JST/JST5.html}。
\end{enumerate}

\begin{J2M}[エントリーリストから 3 つの CSV を作る]
\textbf{練習会モード}を使います。スタート時刻の割り当てをせず，入力順にクラスごとに出力します。

\begin{enumerate}[leftmargin=*,itemsep=0pt]
\item メニューで\textbf{【新規作成】}→ Step 0 でエントリーリスト（CSV／XLSX）をアップロード。
\item 大会設定で\textbf{【練習会モード】にチェック}。
      ゼッケンが要らなければ\textbf{【ゼッケン番号を生成する】を外す}。
\item Step 1 でクラスを確認したら，そのまま\textbf{【生成】}。
      レーン割・シャッフル・制約設定（Step 2・3）は練習会モードでは出ません。
\item Done 画面で ZIP を落とす。\texttt{Startlist.csv} を Mulka2 にドロップし，
      \texttt{RoleStartlists\textbackslash{}Rescue\_ByKana} を救護用に印刷する。
\end{enumerate}

サンプル：\ \url{https://github.com/AtsushiYanaigsawa768/JOY2Mulka/tree/main/sample}
\end{J2M}

\newpage

\Head{当日の手順}

\begin{center}
\small
\begin{tabular}{@{}P{16mm}P{50mm}P{106mm}@{}}
\toprule
\hd{いつ} & \hd{やること} & \hd{中身} \\
\midrule
開場前 & 立ち上げ &
  読み取り機を\textbf{先に}つないでから Mulka2 を起動。
  メインウィンドウ →【EMIT/SI】→【E カード読み取り】で接続。
  COM 番号は後ろに \texttt{USB} などが付いているものを選ぶ \\
開場前 & ポ確カードを読む &
  設置者が持ち帰ったポ確カードを読み，
  \textbf{全コントロールが正しい番号で動いている}ことを確認。
  \tSIAC\ SIAC を使う日は，タッチフリー用の設定になっているステーションかも見る \\
受付中 & 当日申込を入れる &
  メインウィンドウ →【入力】→【当日申込】。
  \textbf{カード番号・氏名・クラス}だけ入れればよい。
  カード備考に\textbf{「レンタル」}と入れておくと回収漏れが減る \\
競技中 & カードを読む &
  読み取り画面を開いたままにする。
  \tEMIT\ E カードは\textbf{読むとカードの時計が止まる}ので，読む前に人を待たせない。
  \tSI\ SI はステーション側に時刻が入っているので順番は問わない \\
競技中 & ペナをさばく &
  MP（ミスパンチ）が出たら本人に説明して確定。
  \textbf{読めなかったカードは横に置き，あとでバックアップから復元} \\
閉鎖前 & 未帰還を確認 &
  読み取り済みの人に対応表の印を付け，\textbf{印の付いていない人を読み上げる}。
  残っている人がいれば救護・運営に伝える \\
撤収後 & 結果を出す &
  メインウィンドウ →【リザルト】で確定。
  ラップと成績はクラウド／ラップセンターで公開する \\
\bottomrule
\end{tabular}
\end{center}

\Head{困ったときの初動}

\begin{center}
\small
\begin{tabular}{@{}P{44mm}P{128mm}@{}}
\toprule
\hd{起きたこと} & \hd{最初にやること} \\
\midrule
読み取り機を認識しない &
  (1) 別の USB ポート　(2) 別のケーブル　(3) ドライバ再インストール　(4) 予備機に交換。
  \textbf{この順で試す}。ケーブルの断線が意外に多い \\
カードが読めない &
  そのカードを横に置き，\textbf{氏名とカード番号を紙に控える}。
  あとでステーションのバックアップメモリから復元する。
  \tEMIT\ E カードならバックアップラベルの穴で通過を証明できる \\
未登録のカードが来た &
  読み取り画面に「未登録」と出る。
  \textbf{その場でカード番号と氏名を控え}，当日申込として入れ直す \\
コースが違った &
  【コース設定変更】で直す。
  すでに読んだ人の記録は\textbf{再計算される}ので，読み直しは不要 \\
PC が固まった &
  \textbf{慌てて電源を切らない}。読み取り済みの記録は \texttt{Changes.dat} に追記済みなので，
  再起動すれば残っている \\
\bottomrule
\end{tabular}
\end{center}

\begin{Fail}
\textbf{起きること：}カードを読んだのに記録が出ない。\\
\textbf{対処：}その人が\textbf{参加者データに入っていない}ことがほとんどです。
当日申込で入れ忘れていないか，クラス名の表記がコースと合っているかを見ます。
\end{Fail}

\begin{Note}[練習会でも守ること]
\begin{itemize}[leftmargin=*,itemsep=0pt]
\item \textbf{レンタルカードは番号順に並べて渡し，番号順に回収する。}
      練習会の紛失はほぼこれで防げます。
\item \textbf{イベントフォルダは帰る前に zip で固めて 1 か所に置く。}
      あとから「あの日の記録は？」と聞かれます。
\item \textbf{読めなかったカード・止めた処理は，紙にその場で書く。}
      記憶に頼ると翌日には消えています。
\end{itemize}
\end{Note}

\vspace{2mm}
\noindent\textcolor{rulec}{\rule{\textwidth}{0.4pt}}\par
\vspace{1mm}
{\footnotesize
入手先：\\
Mulka2 \url{https://mulka2.com/mulka2/ja/index.php/Mulka2}\\
ラップセンター \url{https://mulka2.com/lapcenter/}\\
JOY2Mulka \url{https://github.com/AtsushiYanaigsawa768/JOY2Mulka}\par
この例に登場する人名・所属・カード番号・時刻は架空のものです。
画面上のメニュー名はバージョンによって異なります。}

\end{document}
